\begin{theorem}[严格化问题的双层分裂：测度层刚性与演算层仿射规范]\label{thm:conclusion-strictification-bilayer-measure-affine-gauge}
\leanverified{paper\_conclusion\_strictification\_bilayer\_measure\_affine\_gauge}
设
\[
X\xrightarrow{f}Y\xrightarrow{g}Z
\]
为有限集上的满射，并固定与该局部比较相联系的有限 PW 切片 \(K\) 及其异常 \(2\)-上链
\[
a\in C^2(K;A_G).
\]
则 strictification 问题严格分裂为两个互不混同的层次：
\[
\text{测度层 strictification}
\qquad+\qquad
\text{演算层 strictification}.
\]
其中：
\begin{enumerate}
  \item 测度层没有任何规范余量；唯一可行的 strictification 核正是离散雅可比修正核
  \[
  \widetilde L_{g,f}:\Delta(Z)\to\Delta(Y),
  \qquad
  L_f\widetilde L_{g,f}=L_{g\circ f}.
  \]
  \item 演算层的 strictification 若可行，则其全部解恰构成一个仿射空间
  \[
  \eta_0+Z^1(K;A_G),
  \]
  其中 \(\eta_0\in C^1(K;A_G)\) 满足
  \[
  a=\delta\eta_0.
  \]
\end{enumerate}
换言之，严格化的概率部分被唯一强制，而全部剩余自由度都被压缩为 \(1\)-上链的仿射规范。
\end{theorem}
\begin{proof}
测度层唯一性正是定理 \ref{thm:conclusion-discrete-jacobian-strictification-mediated-defect} 的第一、二部分：凡满足
\[
L_fK=L_{g\circ f}
\]
的 Markov 核 \(K\) 必唯一等于 \(\widetilde L_{g,f}\)。故测度层不存在任何规范余量。

另一方面，正文关于 counterterm strictification 的充要判据已经给出：PW 切片上的 strictification 存在，当且仅当存在 \(1\)-上链 \(\eta\) 满足
\[
a=\delta\eta.
\]
一旦 \(\eta_0\) 为一组解，则另一组 \(\eta\) 与之给出同一 strictification，当且仅当
\[
\delta(\eta-\eta_0)=0,
\]
亦即 \(\eta-\eta_0\in Z^1(K;A_G)\)。故全部解正是仿射空间 \(\eta_0+Z^1(K;A_G)\)。证毕。
\end{proof}

\begin{proposition}[纯标量补偿的充要平坦判据]\label{prop:conclusion-strictification-pure-scalar-compensation-flatness}
\leanverified{paper\_conclusion\_strictification\_pure\_scalar\_compensation\_flatness}
设
\[
X\xrightarrow{f}Y\xrightarrow{g}Z
\]
为有限集上的满射，并令
\[
(L_g\delta_z)(y)=\frac{\mathbf 1_{g(y)=z}}{d_g(z)}
\]
为沿 \(g\)-纤维的朴素均匀提升。则
\[
L_fL_g=L_{g\circ f}
\]
成立，当且仅当 \(d_f\) 在每条 \(g\)-纤维上为常数；等价地，
\[
\widetilde L_{g,f}=L_g.
\]
因此，只要上游纤维基数在某条下游纤维内不平坦，任何 strictification 都不可能由“只依赖 \(z\) 而不区分 \(y\)”的纯标量补偿完成。
\end{proposition}
\begin{proof}
由定理 \ref{thm:conclusion-discrete-jacobian-strictification-mediated-defect}，
\[
(\widetilde L_{g,f}\delta_z)(y)
=
\frac{\mathbf 1_{g(y)=z}\,d_f(y)}{d_{g\circ f}(z)}.
\]
若 \(L_fL_g=L_{g\circ f}\)，则满足 strictification 方程的核必须唯一等于 \(\widetilde L_{g,f}\)，故对每个 \(z\) 与每个 \(y\in g^{-1}(z)\) 都有
\[
\frac1{d_g(z)}=\frac{d_f(y)}{d_{g\circ f}(z)}.
\]
于是 \(d_f(y)\) 在 \(g^{-1}(z)\) 上为常数。

反之，若 \(d_f\) 在每条 \(g\)-纤维上为常数，则上式成立，故 \(\widetilde L_{g,f}=L_g\)，从而
\[
L_fL_g=L_f\widetilde L_{g,f}=L_{g\circ f}.
\]
证毕。
\end{proof}

\begin{theorem}[有限切片全局 strictification 的唯一障碍类]\label{thm:conclusion-finite-slice-strictification-unique-obstruction-class}
设 \(K\) 为任意有限复杂度的 PW 切片，并在每个基本比较方块上插入唯一的测度层 strictification 核 \(\widetilde L_{g,f}\)。则存在一组仅作用于有限部/常数层的 counterterm 门，把整个切片提升为严格合流系统，当且仅当
\[
[a]=0\in H^2(K;A_G).
\]
若解存在，则全部 global strictification 的集合构成一个以
\[
Z^1(K;A_G)
\]
为平移空间的主齐性空间。
\end{theorem}
\begin{proof}
由定理 \ref{thm:conclusion-strictification-bilayer-measure-affine-gauge}，测度层 strictification 已被唯一核 \(\widetilde L_{g,f}\) 完全穷尽；因此在有限切片 \(K\) 上，所有剩余缺陷都集中到异常 \(2\)-上链 \(a\)。

正文的 counterterm strictification 判据表明：存在 global counterterm 系统当且仅当
\[
a=\delta\eta
\]
对某个 \(\eta\in C^1(K;A_G)\) 成立，即当且仅当
\[
[a]=0\in H^2(K;A_G).
\]
一旦 \(\eta_0\) 为一组解，则由定理 \ref{thm:conclusion-strictification-bilayer-measure-affine-gauge}，全部解恰为
\[
\eta_0+Z^1(K;A_G),
\]
故解空间是一条以 \(Z^1(K;A_G)\) 为平移空间的主齐性空间。证毕。
\end{proof}
\leanverified{paper\_conclusion\_finite\_slice\_strictification\_unique\_obstruction\_class}

\begin{theorem}[隐藏子群决定异常最小字母表]\label{thm:conclusion-hidden-subgroup-determines-anomaly-minimal-alphabet}
固定分辨率 \(m\)。称一个有限字母表异常读出器是角色忠实的，若它能够区分全部潜在异常方向，等价地，能够区分 Pontryagin 对偶群
\[
\widehat{N_m^{\mathrm{BC}}}.
\]
则任何此类读出器的输出字母表基数至少为
\[
|N_m^{\mathrm{BC}}|,
\]
因而任何二进制实现都必须满足最小位预算下界
\[
k_{\min}\ge \left\lceil \log_2|N_m^{\mathrm{BC}}|\right\rceil.
\]
\end{theorem}
\begin{proof}
正文已给出：全部潜在异常方向与有限阿贝尔群 \(\widehat{N_m^{\mathrm{BC}}}\) 自然同构。对有限阿贝尔群有
\[
|\widehat{N_m^{\mathrm{BC}}}|=|N_m^{\mathrm{BC}}|.
\]
若某个读出器的字母表基数严格小于 \(|N_m^{\mathrm{BC}}|\)，则由鸽巢原理，\(\widehat{N_m^{\mathrm{BC}}}\) 中至少有两个不同方向会被压缩到同一输出符号，故该读出器不可能角色忠实。

因此角色忠实性强制
\[
|\Sigma_{\mathrm{out}}|\ge |N_m^{\mathrm{BC}}|.
\]
若输出用 \(k\) 个二进制位编码，则 \(|\Sigma_{\mathrm{out}}|\le 2^k\)，故
\[
2^k\ge |N_m^{\mathrm{BC}}|,
\]
从而
\[
k\ge \left\lceil \log_2|N_m^{\mathrm{BC}}|\right\rceil.
\]
证毕。
\end{proof}
